\part{Documentation des données et traitements effectués}
\label{part:doc}
\section{Introduction}

   L'architecture AIDA est conçue pour traiter la metadonnée
conjointement à la donnée : 
en-tête de fichier vecteur contenant des attributs descriptifs,
méta-données sur les acquisitions regroupées dans fichiers \path{.inf}

    Pour être utilisée, affichée, partagée, cette information
doit être sélectionnée, parcourue, mise en forme. C'est le rôle 
des programmes décrits ici.


\section{mk\_doctrt}
     Qu'il soit une mesure brute ou le résultat d'un
calcul, chaque fichier-vecteur contient dans son en-tête 
des informations sur les conditions de sa production. 

Le programme \texttt{mk\_doctrt} produit un document \LaTeX
décrivant une liste de vecteurs, typiquement le contenu d'un
fichier produit et diffusé.

\texttt{mk\_doctrt} a deux syntaxes d'appel :
\begin{itemize}
\item La première forme consiste en une simple liste d'arguments, qui
doivent être en nombre fixe (3) et dans l'ordre : \lstinline[basicstyle=\ttfamily]{mk_doctrt Acq_id  list_vect AA-EXP}
\item La seconde forme utilise le système usuel des lignes
de commandes, ce qui permet des arguments optionnel, et un ordre qui
n'est pas figé :\\
 \lstinline[basicstyle=\ttfamily]
  {mk_doctrt -a Acq_id -f list_vect -d AA-EXP [-l LL] [-u s|m|h|d|F] }
\end{itemize}

\begin{description}
 \item[Acq\_id] identificateur d'acquisition (6 à 8 caractères). Ex : \texttt{me9814}
 \item[list\_vect] fichier contenant la liste des vecteurs à documenter. Ce
               fichier peut être un fichier de type \texttt{.ofs}.  
               Ce nom sera utilisé comme nom du "dataset".
 \item[AA-EXP] répertoire de campagne contenant la doc dans "notes" (ex : 98-FETCH)
 \item[LL] Code ISO-639-1 de la langue (fr pour français, en pour anglais)
 \item[s|m|h|d|F] unité de la fréquence : F pour Hertz, s pour période en seconde, ...
\end{description}

\section{tab\_camp\_HTML\_v2}

Ce script produit un tableau récapitulatif d'une campagne au format HTML.

Il trouve, parmi les fichiers .inf d'une plate-forme, ceux qui correspondent
à une campagne donnée, et en extrait les lieux et heures
de départ et d'arrivée, ainsi que le thème.
